% Clase 6 --- Superficie cerrada que NO contiene la carga (§2.3).
% Un cono con vértice en Q recorta dos casquetes, S2 en la esfera de radio a y
% S1 en la de radio b; las caras laterales tienen n perpendicular a E. Como los
% casquetes son homotéticos, S1/b^2 = S2/a^2 y los flujos se cancelan.
\begin{center}
\begin{tikzpicture}[scale=1]

  % esferas de referencia
  \draw[figaux] (0,0) circle (1.5);
  \draw[figaux] (0,0) circle (3.0);

  % carga en el vértice
  \node[figpos, minimum size=8pt] at (0,0) {};
  \node[figlbl, anchor=north east] at (-0.1,-0.1) {$Q$};

  % cono (bien abierto para que quepan los rótulos)
  \draw[figline] (0,0) -- (45:3.0);
  \draw[figline] (0,0) -- (-5:3.0);

  % casquetes
  \draw[figred, line width=2.4pt] (-5:1.5) arc (-5:45:1.5);
  \draw[figred, line width=2.4pt] (-5:3.0) arc (-5:45:3.0);
  \node[figlbl, figred, anchor=east] at (20:1.35) {$S_2$};
  \node[figlbl, figred, anchor=west] at (20:3.15) {$S_1$};

  % normales salientes de cada casquete (una hacia adentro, la otra hacia afuera)
  \draw[figvec] (8:1.5) -- (8:0.95);
  \node[figlblaux, anchor=north] at (8:0.85) {$\hat n_2$};
  \draw[figvec] (34:3.0) -- (34:3.6);
  \node[figlblaux, anchor=south west] at (34:3.62) {$\hat n_1$};

  % campo radial atravesando el cascarón
  \draw[figvecres] (20:1.7) -- (20:2.8);
  \node[figlbl, figred, anchor=south] at (23:2.3) {$\vec E$};

  % lateral: n perpendicular a E
  \draw[figvec, figgray] (45:2.2) -- ++(135:0.75);
  \node[figlblaux, anchor=south east] at (63:2.75) {$\hat n \perp \vec E$};

  % radios
  \draw[figvecaux] (0,0) -- (-55:1.5);
  \node[figlblaux, anchor=north east] at (-55:1.05) {$a$};
  \draw[figvecaux] (0,0) -- (-100:3.0);
  \node[figlblaux, anchor=east] at (-100:1.9) {$b$};

\end{tikzpicture}\\[0.5em]
{\small\itshape El casquete lejano recibe un campo $b^{2}/a^{2}$ veces menor,
pero su área es $b^{2}/a^{2}$ veces mayor: los dos flujos son iguales y de signo
opuesto, porque las normales salientes apuntan al revés. Toda línea que entra por
$S_2$ sale por $S_1$. El cociente $S/r^{2}$ que aparece acá es el \textbf{ángulo
sólido}.}
\end{center}
